\chapterdivider{11}{Guided application workflows}
  {How should a newcomer plan, run, inspect, and validate a JaxCont study?}
  {Featured applications: six reproducible paths through the example gallery}

\begin{frame}{A research study is a six-stage loop}
\centering
\begin{tikzpicture}[>=Latex,node distance=4.5mm and 4mm,
 stage/.style={rounded corners,draw=JaxBlue,thick,fill=SoftBlue,
   minimum width=26mm,minimum height=11mm,align=center,font=\small},
 inspect/.style={rounded corners,draw=JaxOrange,thick,fill=JaxOrange!10,
   minimum width=38mm,minimum height=20mm,align=center,font=\scriptsize}]
\node[stage] (q) {scientific\\question};
\node[stage,right=of q] (model) {model or\\residual};
\node[stage,right=of model] (seed) {trustworthy\\seed};
\node[stage,below=13mm of seed] (small) {small\\continuation};
\node[inspect,left=of small] (inspect) {inspect / visualize\\[1mm]
  re-evaluate $F$; \texttt{n\_valid}/mask; diagnose stop\\
  configured step bounds; stability; event diagnostics};
\node[stage,left=of inspect] (repeat) {repeat and\\validate};
\draw[->,thick] (q)--(model);
\draw[->,thick] (model)--(seed);
\draw[->,thick] (seed)--(small);
\draw[->,thick] (small)--(inspect);
\draw[->,thick] (inspect)--(repeat);
\draw[->,thick,JaxTeal] (repeat.west) -- ++(-5mm,0) |- (q.west);
\end{tikzpicture}

\vfill
\begin{block}{Start small, then earn complexity}
First establish one known solution and a short, diagnosable branch. Expand the
parameter range, mesh, events, and validation only after that calculation is trustworthy.
\end{block}
\end{frame}

\begin{frame}{Demo cards 1--2: an equilibrium branch and its stability}
\scriptsize
\begin{block}{1 \quad README saddle-node quick start}
\textbf{Objective:} follow $u^2+p=0$ through its fold.\quad
\textbf{File:} \texttt{README.md}, ``Quick start''.\\
\textbf{Run:} \texttt{python -}, paste that Python block, then send EOF.\\
\textbf{Expected:} \texttt{[('fold', 0.0)]} and an interactive annotated branch figure
(no file is saved); a
verification run stored 201 points and reached $\max p\approx-1.19\times10^{-3}$.\\
\textbf{Caution:} no real equilibrium exists for $p>0$; the requested
\texttt{p\_span} endpoint is not evidence that the branch can reach it.
\end{block}

\begin{block}{2 \quad Van der Pol equilibrium Hopf crossing}
\textbf{Objective:} track the origin's stability through $\mu=0$.\quad
\textbf{File:} \texttt{examples/example\_03\_van\_der\_pol.py}.\\
\textbf{Run:} \texttt{cd examples \&\& MPLBACKEND=Agg}\,$\backslash$\\
\hspace*{8mm}\texttt{python example\_03\_van\_der\_pol.py}\\
\textbf{Expected:} \texttt{van\_der\_pol.png}; 83 points, one Hopf event at
$\mu=-1.24\times10^{-10}$, and zero equilibrium residual.\\
\textbf{Caution:} $l_1=0$ here: this is a degenerate center crossing, not a
generic local birth of a small-amplitude cycle.
\end{block}
\end{frame}

\begin{frame}{Demo cards 3--4: cycles and phase-plane geometry}
\scriptsize
\begin{block}{3 \quad Van der Pol periodic orbit}
\textbf{Objective:} refine a simulated cycle, then continue it from $\mu=1$ to $4$.\quad
\textbf{File:} \texttt{examples/example\_10\_van\_der\_pol\_limit\_cycle.py}.\\
\textbf{Run:} \texttt{cd examples \&\& MPLBACKEND=Agg}\,$\backslash$\\
\hspace*{8mm}\texttt{python example\_10\_van\_der\_pol\_limit\_cycle.py}\\
\textbf{Expected:} \texttt{images/example\_10\_van\_der\_pol\_limit\_cycle.png};
138 stable points, $\mu\in[1.000,4.021]$, period $6.6633\to10.2539$.\\
\textbf{Caution:} JaxCont refines the supplied, time-rebased trajectory; it
does not simulate to discover the initial cycle.
\end{block}

\begin{block}{4 \quad FitzHugh--Nagumo phase plane}
\textbf{Objective:} connect a Hopf marker on the branch to 2D flow geometry.\quad
\textbf{File:} \texttt{examples/example\_12\_fitzhugh\_nagumo\_phase\_plane.py}.\\
\textbf{Run:} \texttt{cd examples \&\&}\,$\backslash$\\
\hspace*{8mm}\texttt{python example\_12\_fitzhugh\_nagumo\_phase\_plane.py}\\
\textbf{Expected:} stdout \texttt{hopf at I = 0.3313} and two displayed
interactive figures (the script does not save them).\\
\textbf{Caution:} the phase-plane API is for two-state autonomous systems at
one frozen parameter; its trajectories are simulations, not continued cycles.
\end{block}
\end{frame}

\begin{frame}{Demo cards 5--6: phase sensitivity and an independent check}
\scriptsize
\begin{block}{5 \quad Circle-system iPRC and parameter dPRC \quad \mainbadge}
\textbf{Objective:} compute $Z(t)$ and its total parameter derivative.\quad
\textbf{File:} \texttt{examples/example\_13\_phase\_response\_curve.py}.\\
\textbf{Run:} \texttt{cd examples \&\& MPLBACKEND=Agg}\,$\backslash$\\
\hspace*{8mm}\texttt{python example\_13\_phase\_response\_curve.py}\\
\textbf{Expected:} \texttt{images/example\_13\_phase\_response\_curve.png}
with iPRC and $d/d\rho$ panels.\\
\textbf{Caution:} \texttt{dprc\_curve} re-solves $U(p)$; it is not a partial
derivative of \texttt{prc\_curve} at a frozen orbit.
\end{block}

\begin{block}{6 \quad Independent PRC shooting validation \quad \mainbadge}
\textbf{Objective:} compare collocation PRCs with IVP shooting and finite differences.\quad
\textbf{File:} \texttt{examples/example\_14\_prc\_shooting\_validation.py}.\\
\textbf{Run:} \texttt{cd examples \&\& MPLBACKEND=Agg}\,$\backslash$\\
\hspace*{8mm}\texttt{python example\_14\_prc\_shooting\_validation.py}\\
\textbf{Expected:} \texttt{images/example\_14\_prc\_shooting\_validation.png};
max PRC/dPRC discrepancies $3.27\times10^{-5}/4.46\times10^{-5}$ (circle)
and $1.87\times10^{-4}/5.54\times10^{-4}$ (Van der Pol).\\
\textbf{Caution:} this one-off script validates the PRC computation, not
independent periodic-orbit discovery; its orbit seed is still reconverged by JaxCont.
\end{block}
\end{frame}

\begin{frame}{How to read a result}
\centering
\begin{tikzpicture}[>=Latex,node distance=5mm and 8mm,
 q/.style={diamond,aspect=2.4,draw=JaxBlue,thick,fill=SoftBlue,
   align=center,font=\small,inner sep=1.5pt},
 box/.style={rounded corners,draw=JaxTeal,thick,fill=JaxTeal!9,
   text width=51mm,minimum height=13mm,align=left,font=\scriptsize},
 final/.style={rounded corners,draw=JaxOrange,thick,fill=JaxOrange!10,
   text width=49mm,minimum height=11mm,align=left,font=\scriptsize}]
\node[q] (stop) {Did continuation\\stop early?};
\node[box,below left=6.5mm and 3mm of stop] (yes) {Re-evaluate $F(U,p)$ on valid
states; check finite values and endpoint coverage. Review configured
\texttt{ds\_min}/\texttt{max\_steps}; rejected attempts add unreported work.};
\node[box,below right=6.5mm and 3mm of stop] (no) {Inspect stability and event evidence:
eigenvalues or Floquet multipliers, detector diagnostics, and local refinement residuals.};
\node[final,below=6.5mm of yes] (important) {\textbf{Important conclusion:}
rerun with changed step bounds, resolution, and direction.};
\node[final,below=6.5mm of no] (publish) {\textbf{Publication claim:}
add an analytic oracle or an independent method/package.};
\draw[->,thick] (stop)--node[above left,font=\scriptsize]{yes}(yes);
\draw[->,thick] (stop)--node[above right,font=\scriptsize]{no}(no);
\draw[->,thick] (no)--node[left,font=\scriptsize]{important}(important);
\draw[->,thick] (no)--node[right,font=\scriptsize]{publication}(publish);
\end{tikzpicture}

\vfill
\begin{alertblock}{Do not infer a physical branch ending from missing points}
The public result does not store residual norms, accepted step sizes, rejection
history, or a termination reason. Re-evaluate $F$ on valid states and compare
the last valid point with the configured endpoint and stop conditions.
\end{alertblock}
\end{frame}
